\chapter{Proposed Methodology}
\label{ch:methodology}

This chapter describes the system, the data, the degradation simulator, the models and the evaluation protocol. Where an implementation property affects how a result should be interpreted, it is stated in the section it affects and summarised in Section~\ref{sec:audit}.

\section{System overview}
\label{sec:overview}

For regularisation investigations, the framework maintains a compact backbone (MobileNetV3-Small~\cite{howard2019mobilenetv3} in downstream exploratory experiments) and varies the training objective. Waveforms are cropped, converted to normalised log-mel features and classified. During robustness training each crop is also degraded and passed through the same network, and coherence losses tie the two outputs together.

\subsection{Design principles}

\begin{itemize}
  \item \textbf{Same inference graph.} All robustness machinery is used in training. At test time the model computes one mel spectrogram and runs one forward pass.
  \item \textbf{A clean-trained reference.} Mode~M0 (class-weighted cross-entropy on clean input) provides the clean reference baseline.
  \item \textbf{Agreement between views.} A clean clip and a degraded copy of it carry the same label, so their outputs should be similar. Plain augmentation only shows the model degraded data; it does not ask the two outputs to agree.
\end{itemize}

% Master proposed-methodology flowchart (replaces pipeline_overview.png)
\begin{figure}[htbp]
\centering
\resizebox{\textwidth}{!}{%
\begin{tikzpicture}
% (a) data preparation and paired views
\node[blk,text width=7.6cm] (in) at (0,0) {\textbf{AT-ADD Track~2 clip}\\mono, 24\,kHz, label $y\in\{0,1\}$, audio type $t$};
\node[blk,text width=9.2cm] (seg) at (0,-1.6) {\textbf{Segmentation}\\4\,s crop (random in training, centre in evaluation) or repeat-padding, $N=96{,}000$};
\node[blk,text width=4.6cm] (xc) at (-3.8,-3.3) {\textbf{Clean view} $x$};
\node[trn,text width=6.4cm] (g) at (3.7,-3.3) {\textbf{Degradation simulator} $\tilde{x}=g(x;c)$\\MP3 $\cdot$ noise $\cdot$ resampling $\cdot$ telephony;\\severity $\rho\le\rho_{\max}(e)$ (curriculum)};
% (b) frontend and backbone
\node[blk,text width=5.6cm] (mc) at (-3.8,-6.3) {\textbf{Log-mel frontend} $\phi$\\STFT $\to$ 80 mel bins $\to$ log $\to$ z-score\\output $\hat{\mathbf L}_c$};
\node[trn,text width=5.6cm] (md) at (3.8,-6.3) {\textbf{Log-mel frontend} $\phi$ (identical)\\STFT $\to$ 80 mel bins $\to$ log $\to$ z-score\\output $\hat{\mathbf L}_d$};
\node[blk,text width=10.4cm] (bb) at (0,-8.6) {\textbf{Shared MobileNetV3-Small backbone} ($\approx$1.52\,M parameters)\\inverted residual blocks $\to$ feature map $F$ $\to$ global pool $\to$ 2-class head};
\node[blk,text width=5.2cm,minimum height=1.5cm] (oc) at (-3.8,-10.7) {\textbf{Clean outputs}\\logits $z_c$, score $q_c$,\\8-band signature $s_c$};
\node[trn,text width=5.2cm,minimum height=1.5cm] (od) at (3.8,-10.7) {\textbf{Degraded outputs}\\logits $z_d$, score $q_d$,\\8-band signature $s_d$};
% (c) objectives
\node[los,text width=3.2cm,minimum height=1.9cm] (lcls) at (-5.8,-14.0) {$\mathcal{L}_{\mathrm{cls}}$\\dual weighted CE\\$\mathcal{L}_{\mathrm{CE}}(z_c,y)+\beta\mathcal{L}_{\mathrm{CE}}(z_d,y)$};
\node[los,text width=3.2cm,minimum height=1.9cm] (lpc) at (-1.95,-14.0) {$\mathcal{L}_{\mathrm{PC}}$\\prediction coherence\\$(q_c-q_d)^2$};
\node[los,text width=3.2cm,minimum height=1.9cm] (lft) at (1.95,-14.0) {$\mathcal{L}_{\mathrm{FTSC}}$\\signature coherence\\$1-\cos(s_c,s_d)$ per band};
\node[los,text width=3.2cm,minimum height=1.9cm] (gate) at (5.8,-14.0) {\textbf{Robustness gate}\\$w_i=\alpha_{t_i}\kappa_i\eta(\rho_i)$\\type, confidence, severity};
\node[los,text width=12.6cm] (tot) at (0,-16.5) {\textbf{Total objective (mode M3)}\\$\mathcal{L}_{\mathrm{M3}}=\mathcal{L}_{\mathrm{cls}}+\bar{w}\,(\gamma_{\mathrm{pc}}\mathcal{L}_{\mathrm{PC}}+\gamma_{\mathrm{ftsc}}\mathcal{L}_{\mathrm{FTSC}})$};
% (d) optimisation and (e) inference side by side
\node[blk,text width=6.8cm] (opt) at (-3.9,-19.6) {\textbf{Optimization}\\AdamW, lr $10^{-3}$, weight decay $10^{-4}$, batch 32, $\le$20 epochs; gradients update the shared backbone};
\node[blk,text width=6.8cm] (sel) at (-3.9,-21.6) {\textbf{Robustness-aware checkpoint selection}\\$J=\mathrm{Macro\text{-}F1}_{\mathrm{dev}}-\lambda\,\widehat{\Delta\mathrm{F1}}_{\mathrm{dev}}$, $\lambda=0.5$};
\node[dep,text width=6.8cm] (inf) at (3.9,-19.6) {\textbf{Deployed detector} (inference graph unchanged)\\one log-mel pass $\to$ one forward pass\\$q=p_1$, $\hat y=\mathbb 1[q\ge0.5]$};
\node[dep,text width=6.8cm] (ev) at (3.9,-21.6) {\textbf{Evaluation} on held-out clips\\Macro-F1, accuracy, AUROC, EER,\\AT-ADD Macro-F1};
% arrows
\draw[arr] (in)--(seg);
\draw[arr] ([xshift=-1.6cm]seg.south)--(xc.north);
\draw[arr] ([xshift=1.6cm]seg.south)--(g.north);
\draw[arr] (xc)--(mc);
\draw[arr] (g)--(md);
\draw[arr] (mc.south)--(mc.south |- bb.north);
\draw[arr] (md.south)--(md.south |- bb.north);
\draw[arr] (oc.north)--(oc.north |- bb.south);
\draw[arr] (od.north)--(od.north |- bb.south);
\draw[thick] (oc.south)--(oc.south |- 0,-12.1);
\draw[thick] (od.south)--(od.south |- 0,-12.1);
\draw[thick] (-5.8,-12.1)--(5.8,-12.1);
\foreach \n in {lcls,lpc,lft,gate}{\draw[arr] (\n.north |- 0,-12.1)--(\n.north);}
\foreach \n in {lcls,lpc,lft,gate}{\draw[arr] (\n.south)--(\n.south |- tot.north);}
\draw[arr] (tot.south -| opt.north)--(opt.north);
\draw[arr] (opt)--(sel);
\draw[arr] (sel.east)--++(0.45,0)|-node[pos=0.25,left,font=\scriptsize,align=right]{best\\checkpoint}(inf.west);
\draw[arr] (inf)--(ev);
% stage frames
\begin{scope}[on background layer]
\node[stage,fit=(in)(seg)(xc)(g)] (sA) {};
\node[stage,fit=(mc)(md)(bb)(oc)(od)] (sB) {};
\node[stage,fit=(lcls)(lpc)(lft)(gate)(tot)] (sC) {};
\node[stage,fit=(opt)(sel)] (sD) {};
\node[stage,fit=(inf)(ev),draw=deployedge] (sE) {};
\end{scope}
\node[stlab,fill=white,anchor=center] at (sA.north) {(a) Data preparation and paired views};
\node[stlab,fill=white,anchor=center] at (sB.north) {(b) Shared frontend and backbone};
\node[stlab,fill=white,anchor=center] at ([xshift=-4.6cm]sC.north) {(c) FT-RCL training objectives};
\node[stlab,fill=white,anchor=west] at ([xshift=0.3cm]sD.north west) {(d) Optimization};
\node[stlab,fill=white,anchor=center] at (sE.north) {(e) Inference and evaluation};
\end{tikzpicture}}
\caption{Overall end-to-end workflow of the proposed methodology. Orange blocks are used only during training; the deployed detector (green) keeps the single-view baseline inference graph.}
\label{fig:pipeline-overview}
\end{figure}


\section{Problem formulation}
\label{sec:formulation}

Let $x \in \mathbb{R}^{N}$ be a mono waveform with sampling rate $f_s = 24{,}000$\,Hz and fixed length $N = 4 f_s = 96{,}000$ samples (four seconds). The label $y \in \{0,1\}$ indicates real ($0$) or fake ($1$). The classifier produces logits $z \in \mathbb{R}^2$ and probabilities $p = \mathrm{softmax}(z)$. The \textbf{detection score} used for ROC analysis and EER is the fake-class probability $q = p_1$.

For robustness training, a condition $c$ from a finite catalogue defines a waveform-domain operator $g(\cdot\,; c)$. The paired sample $(x, \tilde{x})$ satisfies $\tilde{x} = g(x; c)$ with the same label $y$ and the same temporal crop. Mel extraction $\phi(\cdot)$ is applied \textbf{identically} to both waveforms:
\begin{equation}
\mathbf{X} = \phi(x), \qquad \tilde{\mathbf{X}} = \phi(\tilde{x}).
\label{eq:pair}
\end{equation}
Degradation never alters the mel hyperparameters or normalization statistics, which ensures a fair comparison with the clean baseline.

\section{Dataset and evaluation protocol}
\label{sec:data}

\subsection{AT-ADD Track 2}
\label{sec:data-track2}

AT-ADD Track 2 provides real/fake labels and audio-type annotations for speech, singing, music, and environmental sound. Generator metadata is available for fake audio samples. The present study uses project-specific manifest files prepared from the broader AT-ADD Track 2 dataset. The sample counts reported in the following sections refer to these manifests and should not be interpreted as the total size of the complete benchmark.

The three manifests define the working dataset used in this study: 60,000 training samples, 22,536 development samples, and 44,289 evaluation samples. All records listed in these manifests were used in the reported experiments. The partitioning procedure and the resulting generator and speaker separation are described in Section~\ref{sec:partitioning}, while Section~\ref{sec:data-used} specifies how the manifest-defined data were used for training, validation, and evaluation.

\subsection{Dataset Partitioning and Split Integrity}
\label{sec:partitioning}
\label{sec:overlap}

The original training and development data can share fake-audio generators, which may lead to optimistic validation results when a detector learns generator-specific characteristics. To address this issue, the project-specific data were organised into generator-disjoint training, development, and evaluation manifests. The resulting partitions were audited to examine generator overlap and speaker-ID overlap.

The three manifests contain 126,825 audio samples in total, distributed across 60,000 training samples, 22,536 development samples, and 44,289 evaluation samples. These counts describe the working dataset used in this study rather than the complete AT-ADD Track 2 collection. Table~\ref{tab:dataset-manifests} summarises the distribution of samples across the three manifests by class and audio type.

\begin{table}[htbp]
\centering
\caption{Sample distribution across the manifest-defined dataset partitions.}
\label{tab:dataset-manifests}
\label{tab:dataset-eval}
\small
\begin{tabular}{@{}lrrrrrrr@{}}
\toprule
Partition & Real & Fake & Total & Speech & Singing & Music & Sound \\
\midrule
Training    & 16{,}000 & 44{,}000 & 60{,}000  & 34{,}000 & 8{,}000 & 10{,}000 & 8{,}000 \\
Development & 6{,}536  & 16{,}000 & 22{,}536  & 12{,}000 & 4{,}000 & 2{,}536  & 4{,}000 \\
Evaluation  & 6{,}536  & 37{,}753 & 44{,}289  & 27{,}000 & 6{,}753 & 4{,}536  & 6{,}000 \\
\midrule
\textbf{Total} & \textbf{29{,}072} & \textbf{97{,}753} & \textbf{126{,}825} & \textbf{73{,}000} & \textbf{18{,}753} & \textbf{17{,}072} & \textbf{18{,}000} \\
\bottomrule
\end{tabular}
\end{table}

The manifest-defined dataset is imbalanced towards fake audio. Fake samples account for approximately 73.33\% of the training set, 71.00\% of the development set, and 85.24\% of the evaluation set. This imbalance is important when interpreting pooled evaluation metrics, particularly because the audio-type distributions also differ across partitions. Accordingly, pooled metrics are considered alongside audio-type-averaged Macro-F1 to assess performance across speech, singing, music, and sound.

\paragraph{Generator-wise partitioning.} The generator metadata was examined to determine whether fake-audio generators were shared across the three manifests. The training set contains 21 distinct generator names, the development set contains 8, and the evaluation set contains 38. No exact generator name is shared between any pair of partitions. This means that the fake-audio generator groups in the training, development, and evaluation manifests are disjoint according to the generator labels recorded in the files.

The training manifest contains 44,000 fake samples from 21 generators. AudioLDM2 contributes 4,000 samples. AudioLDM, DiffSpeech, FastDiff, FastPitch, FastSpeech2, Glow-TTS, GradTTS, Index-TTS, Kokoro, MelGAN, MusicLDM, NeuCoSVC\_inhouse, Parler-TTS (mini), PortaSpeech (normal), ProDiff, So-VITS-SVC\_inhouse, Stable Audio open, StarGANv2-VC, Tacotron2, and WaveGlow each contribute 2,000 samples.

The development manifest contains 16,000 fake samples from eight generators: Audiogen, CosyVoice, Llasa1B, MeloTTS, MusicGen, SeedVC, Tortoise-TTS, and WaveNet. Each contributes 2,000 samples.

The evaluation manifest contains 37,753 fake samples from 38 generator labels. RVC-FSD contributes 753 samples, while ACE-Step, Auffusion, BigVGAN, Both-Ears-Wide-Open, ChatTTS, CosyVoice2, CosyVoice3, DiffGANTTS, E2TTS, F5TTS, FireRedTTS2, Fish, Free-SVC, GPTSoVITS, HiFiGAN, IndexTTS1.5, IndexTTS2, KNN-SVC, Llasa3B, Llasa8B, MBMelGAN, MeanAudio, Musictango, OpenVoice2, ParallelWaveGAN, SSM-TTM, SeedVC1, So-VITS-SVC-FSD, SparkTTS, StepAudioTTS, StyleMelGAN, StyleSpeech, StyleTTS2, UniAudio2, Uniaudio2, VEVO1.5, and VITS each contribute 1,000 samples. Generator names are retained as recorded in the manifests.

\paragraph{Speaker-ID overlap audit.} Speaker metadata was also examined to identify potential overlap between the partitions. After excluding missing identifiers, represented by \texttt{-}, and the generic \texttt{single-speaker} marker, the manifests contain 1,668 distinct explicit speaker IDs in training, 590 in development, and 1,097 in evaluation. No common explicit speaker IDs were found between training and development, training and evaluation, or development and evaluation.

Speaker identifiers are not available for all samples. The \texttt{speaker\_id} field is missing for 52,000 training records, 18,536 development records, and 40,289 evaluation records. Consequently, the audit establishes non-overlap among the available explicit speaker identifiers but cannot independently establish speaker-level separation for recordings without identifiable speaker metadata. Generator disjointness and speaker-ID disjointness are therefore treated as separate properties of the dataset partitioning.

Overall, the manifest audit establishes disjoint generator labels and disjoint explicit speaker identifiers across the three partitions. These checks support evaluation of generalisation to unseen generators and speakers where the metadata permits verification. The conclusions apply to the manifest-defined working dataset and should not be extended to samples outside these manifests without a corresponding audit.

\subsection{Data Used in the Reported Experiments}
\label{sec:data-used}

The reported experiments use all records contained in the three project-specific manifests. The baseline, FT-RCL, and STAR-Edge models use the 60,000 samples in the training manifest for model fitting and the 22,536 samples in the development manifest for validation and checkpoint selection. Final evaluation is performed on all 44,289 samples in the evaluation manifest.

The training manifest contains 16,000 real and 44,000 fake samples. The development manifest contains 6,536 real and 16,000 fake samples, while the evaluation manifest contains 6,536 real and 37,753 fake samples. All four audio categories are represented in each partition. No additional sampling was performed from these manifests for the reported experiments.

The evaluation results therefore correspond to the complete evaluation manifest rather than smaller exploratory subsets. However, the 44,289 evaluated samples represent the evaluation data selected for this study and should not be described as the entirety of the broader AT-ADD Track 2 collection unless that scope is independently established.

No evaluation-manifest samples are used for parameter updates or checkpoint selection. The development set is used for validation, while the evaluation set is reserved for final performance assessment. Since the generator labels are disjoint across the manifests, the evaluation protocol measures generalisation to generator labels not present in the training or development data. The reported results should be interpreted in the context of the manifest-defined dataset size, its class imbalance, its audio-type distribution, and the limitations of the available speaker metadata.

\section{Audio preprocessing and log-mel frontend}
\label{sec:frontend}

\subsection{Loading and segmentation}

Audio clips are converted to mono and resampled to $f_s = 24\,\text{kHz}$ to match the AT-ADD Track~2 evaluation standard. Each clip is formatted to a fixed duration of $N = 4 \times f_s = 96{,}000$ samples (four seconds). A recording longer than four seconds is cropped to length $N$ using a random start offset during training (providing temporal translation augmentation) and a deterministic centre crop during evaluation. A recording shorter than four seconds ($L < N$) is extended to length $N$ by periodic tiling (repeat-padding), $x_{\mathrm{pad}}[n] = x[n \pmod L]$. As justified in Section~\ref{sec:duration-norm}, repeat-padding is adopted as a controlled preprocessing choice to avoid the deterministic trailing-silence cue introduced by zero-padding, recognizing that repetition boundaries may introduce a waveform discontinuity.

\subsection{Short-time Fourier transform and mel filtering}

With a Hann window $w[n]$, frame length $L_w = 600$ samples (25\,ms), hop $H = 240$ samples (10\,ms), and FFT size $N_{\mathrm{FFT}} = 1024$ (the window is zero-padded to the FFT size),
\begin{equation}
X(m,k) = \sum_{n=0}^{L_w-1} x[n + mH]\, w[n]\, e^{-j2\pi kn / N_{\mathrm{FFT}}}, \qquad P(m,k) = |X(m,k)|^2 .
\label{eq:stft}
\end{equation}
Frames are centred using reflection padding, so a four-second segment yields $T_m = 1 + \lfloor N/H \rfloor = 401$ frames. A mel filter bank $\mathbf{W}_{\mathrm{mel}} \in \mathbb{R}^{80 \times 513}$ with 80 bands (Slaney normalization) gives the mel power $\mathbf{P}_{\mathrm{mel}} = \mathbf{W}_{\mathrm{mel}} \mathbf{P}^{\top}$, and the log-mel features are
\begin{equation}
\mathbf{L} = \log(\mathbf{P}_{\mathrm{mel}} + \varepsilon), \qquad \varepsilon = 10^{-10}.
\label{eq:logmel}
\end{equation}

\subsection{Global normalization}

The mean $\mu$ and standard deviation $\sigma$ are computed across all time-frequency bins of the 60{,}000-sample training manifest ($\mu \approx -9.128$, $\sigma \approx 6.774$) and applied globally to training, development, and evaluation splits:
\begin{equation}
\hat{\mathbf{L}} = (\mathbf{L} - \mu) / \sigma .
\label{eq:norm}
\end{equation}
The normalized tensor $\hat{\mathbf{L}}$ is formatted as a single-channel feature representation of shape $1 \times 80 \times T_m$ for input to the convolutional backbones (Figure~\ref{fig:mel-frontend}).

\begin{figure}[htbp]
\centering
\resizebox{\linewidth}{!}{%
\begin{tikzpicture}
\node[blk,text width=2.1cm] (w) at (0,0) {\textbf{Waveform}};
\node[blk,text width=1.8cm] (s) at (2.5,0) {\textbf{STFT}};
\node[blk,text width=2.0cm] (m) at (5.0,0) {\textbf{Mel bank}};
\node[blk,text width=1.5cm] (l) at (7.5,0) {\textbf{log}};
\node[blk,text width=2.0cm] (z) at (10.0,0) {\textbf{Z-score}};
\draw[arr] (w)--(s);
\draw[arr] (s)--(m);
\draw[arr] (m)--(l);
\draw[arr] (l)--(z);
\end{tikzpicture}}
\caption{Audio preprocessing and log-mel feature extraction.}
\label{fig:mel-frontend}
\end{figure}


\section{Synthetic degradation (paired views)}
\label{sec:degradation}

\subsection{Degradation catalogue}

Degradations are applied to the waveform before $\phi$. Table~\ref{tab:degradations} lists the catalogue used in training. MP3 coding uses FFmpeg (libmp3lame), additive noise is white Gaussian noise scaled to a target SNR, and the telephony condition is a fourth-order Butterworth band-pass. Each degraded clip records its family, severity and random seed.

Two properties of the catalogue need to be stated. First, the resampling condition (24\,kHz $\rightarrow$ 16\,kHz $\rightarrow$ 24\,kHz) does nothing to this corpus, because the audio is natively 16\,kHz; it is an identity operation, which explains why its measured $\Delta M$ is exactly zero (Section~\ref{sec:robustness-pilot}). Second, in the implementation used for the reported runs, the noise of a deterministic degradation was drawn from the same seed for every clip, so all clips received the same noise waveform. This has been corrected: the seed now depends on the condition and the clip identifier.

\begin{table}[htbp]
\centering
\caption{Degradation families and severity levels of the training catalogue.}
\label{tab:degradations}
\small
\begin{tabular}{@{}lll@{}}
\toprule
Family & Setting & Severity $\rho$ \\
\midrule
Resample & 24\,kHz $\rightarrow$ 16\,kHz $\rightarrow$ 24\,kHz (identity on this corpus) & 1 \\
MP3 codec & 128 / 64 / 32\,kbps & 2 / 3 / 4 \\
Additive noise & SNR 20 / 10 / 5\,dB & 2 / 3 / 4 \\
Telephony & Band-pass 300--3400\,Hz & 3 \\
\bottomrule
\end{tabular}
\end{table}

The training catalogue and the evaluation conditions are the same families. A test of generalisation therefore needs conditions that were not used in training, for instance AAC, Opus, G.711 telephony, reverberation and packet loss. Such conditions are implemented in the evaluation code but their results are not available (Section~\ref{sec:robustness-pilot}).

\figimage{spectrogram-pairs}{Clean and degraded versions of one evaluation clip as waveform and log-mel spectrogram.}{figures/plots/btp_report/spectrogram_pairs.png}

\subsection{Severity levels and curriculum}

The severity $\rho$ is an ordinal rank used to schedule the curriculum and to scale the gate. It increases within a family as the bitrate or SNR falls, but the ranks of different families are a design choice and not a calibrated quantity: MP3 at 128\,kbps and noise at 20\,dB share rank~2 without being equally damaging. The agreement between assigned rank and measured drop could not be checked because the degraded-audio results are not valid.

Let $\rho_{\max}(e)$ be the largest severity allowed at epoch $e$, with milestones at epochs 6 and 13:
\begin{equation}
\rho_{\max}(e) = \begin{cases}
1 & e < 6, \\
3 & 6 \le e < 13, \\
4 & e \ge 13 .
\end{cases}
\label{eq:curriculum}
\end{equation}
Only conditions with $\rho \le \rho_{\max}(e)$ should be sampled (Figure~\ref{fig:m3-gating}). Two facts limit what this schedule meant in practice. The FT-RCL run stopped at epoch 12 by early stopping, so it never reached the $\rho=4$ stage, and its best checkpoint is from epoch 7. More importantly, the training loader used persistent worker processes, each of which keeps its own copy of the dataset; the severity cap was set in the main process and never reached the workers, so every worker sampled all severities from the first epoch. The curriculum was therefore not in effect in the reported runs. The cap is now kept in shared memory and each worker has its own random stream.

\figimage{m3-gating}{Intended severity curriculum: maximum allowed severity $\rho_{\max}$ versus epoch. In the reported runs the cap did not reach the data-loading workers.}{figures/plots/btp_report/m3_gating_curriculum.png}

The severity factor of the M3 gate (Section~\ref{sec:losses}) is
\begin{equation}
\eta(\rho) = \min\!\left(1,\ \frac{\rho}{\rho_{\max}(e)}\right).
\label{eq:eta}
\end{equation}

\section{Clean baseline model selection (mode M0)}
\label{sec:baseline}

Three candidate detector configurations were evaluated during the clean baseline stage before selecting the model to be carried forward into subsequent experiments, representing a structured representation progression:
\begin{itemize}
  \item \textbf{E1 (Compact CNN reference):} used a scratch-trained MobileNetV3-Small~\cite{howard2019mobilenetv3} operating on the compact 80-bin log-mel representation, with depthwise-separable convolutions and inverted bottlenecks ($1{,}519{,}618$ parameters), drawing architectural motivation from lightweight vision-backbone studies in audio deepfake detection~\cite{pham2024deepfake}.
  \item \textbf{E2 (Temporal baseline):} extended the MobileNetV3-Small convolutional representation with a one-layer GRU with hidden size 128 ($1{,}198{,}082$ parameters) to investigate whether explicit temporal sequence modelling improved the compact detector over global spatial pooling, following convolutional-recurrent designs evaluated in literature~\cite{pham2024deepfake}.
  \item \textbf{E3 (Strong learned representation):} adapted PANNs MobileNetV2~\cite{kong2020panns} initialized with AudioSet-pretrained weights and a newly initialized binary real/fake classification head ($3{,}537{,}250$ parameters). Motivated by evidence that pretrained audio representations transfer rich discriminative priors~\cite{combei2024wavlm,kong2020panns}, this experiment investigated whether general-purpose audio pretraining provided a stronger starting point for deepfake detection than learning the representation entirely from scratch.
\end{itemize}

All three configurations were evaluated using three independent training seeds (seeds 42, 1337, and 2024) under the standardized baseline protocol with gradient clipping ($\text{MAX\_GRAD\_NORM}=5.0$). E3 achieved the strongest clean performance across all metrics (type-averaged Macro-F1 of $0.6274 \pm 0.0158$, AUROC of $0.8800 \pm 0.0032$, and EER of $0.1993 \pm 0.0028$) and demonstrated substantially lower variation across seeds. It was therefore selected as the clean reference baseline \textbf{M0}. From this point onward, \textbf{M0 refers specifically to E3} (Figure~\ref{fig:mobilenet-baseline}).

\begin{figure}[htbp]
\centering
\resizebox{\linewidth}{!}{%
\begin{tikzpicture}
\node[blk,text width=2.4cm] (in) at (0,0) {$1\times80\times T$\\log-mel};
\node[blk,text width=2.2cm] (bn) at (3.0,0) {\textbf{Fresh BN0}\\$80$ mel bins};
\node[blk,text width=3.4cm] (mn) at (6.6,0) {\textbf{PANNs MobileNetV2}\\AudioSet-pretrained};
\node[blk,text width=2.8cm] (gp) at (10.6,0) {\textbf{Pooling}\\Freq-mean $\to$ Temp max+mean};
\node[blk,text width=2.4cm] (fc) at (14.2,0) {\textbf{Pretrained FC1}\\$1024$ + Dropout};
\node[blk,text width=2.2cm] (hd) at (17.4,0) {\textbf{2-class head}\\binary real/fake};
\draw[arr] (in)--(bn);
\draw[arr] (bn)--(mn);
\draw[arr] (mn)--(gp);
\draw[arr] (gp)--(fc);
\draw[arr] (fc)--(hd);
\end{tikzpicture}}
\caption{Selected clean baseline architecture M0 (PANNs MobileNetV2 with AudioSet-pretrained feature extractor, fresh 80-bin BN0, and binary classification head; $3{,}537{,}250$ parameters).}
\label{fig:mobilenet-baseline}
\label{fig:panns-baseline}
\end{figure}


\subsection{Class-weighted cross-entropy loss}

With class counts $n_0, n_1$ and total $n = n_0 + n_1$, the weights $\omega_j = n / (2 n_j)$ balance the real and fake classes:
\begin{equation}
\mathcal{L}_{\mathrm{CE}}(z, y) = -\sum_{j \in \{0,1\}} \omega_j \,\mathbb{1}[y=j]\,\log p_j .
\label{eq:ce}
\end{equation}
For the $16{,}000:44{,}000$ (real:fake) ratio of the training manifest ($n=60{,}000$), this gives $\omega_0 = 1.875$ for genuine audio and $\omega_1 \approx 0.682$ for synthetic audio. Mode \textbf{M0} (clean-only training with selected model E3) minimizes $\mathcal{L}_{\mathrm{CE}}$ on clean batches only.

\section{FT-RCL detector}
\label{sec:ftrcl}

\subsection{Architecture and subband signatures}

The FT-RCL detector shares the MobileNetV3 backbone and also returns the last feature map $F \in \mathbb{R}^{B \times C \times H \times W}$ with $C=576$ (Figure~\ref{fig:ft-rcl-paired}). The signature module splits the frequency axis into $K=8$ bands. For band $k$ with rows $[h_k, h_{k+1})$,
\begin{equation}
b_k = \mathrm{mean}_{h,w}\big(F_{:,:,h_k:h_{k+1},:}\big) \in \mathbb{R}^{B \times C}.
\label{eq:subband}
\end{equation}
The bands are stacked into $s^{(0)} \in \mathbb{R}^{B \times K \times C}$ and normalised along the channel axis:
\begin{equation}
s = \mathrm{normalize}\big(s^{(0)},\ p=2,\ \mathrm{dim}=-1\big).
\label{eq:signature}
\end{equation}

The question of how eight bands arise from a feature map with few rows has a direct answer: they do not. The backbone reduces frequency resolution by a factor of 32, so an input with 80 mel bins gives $H=3$ rows at the last layer. The implementation used for the reported runs assigned one row to each of the first three bands and left bands four to eight empty. The mean of an empty slice is undefined, so five of the eight signature bands were not-a-number, the logged FTSC and total loss were not-a-number in every epoch (the loss columns of the saved training logs are empty), and the training loss could not be used for monitoring. The gradients were nevertheless finite, because the undefined bands are not connected to the feature map, so training went on and the FTSC term constrained only the first three bands (one row each). We reproduced this behaviour: with the old module the loss is undefined while all parameter gradients are finite. An exploratory conjecture that the loss was unstable under mixed precision was unfounded; the exploratory runs used a CPU and no mixed precision (Chapter~\ref{ch:setup}).

The corrected module averages over time and then divides the $H$ rows into $K$ possibly overlapping bins with the boundaries used by adaptive average pooling, so that no band is empty. With $H=3$ the eight bands are still an interpolated view of three rows and not eight independent frequency bands. Taking the feature map from an earlier stage with finer frequency resolution would be a better design and is left for future work. The results of Chapter~\ref{ch:results} come from the old module, so they describe FTSC on three bands.

\begin{figure}[htbp]
\centering
\resizebox{0.92\linewidth}{!}{%
\begin{tikzpicture}
\node[blk,text width=2.6cm] (c) at (-2.2,1.1) {\textbf{Clean mel}};
\node[trn,text width=2.6cm] (d) at (-2.2,-1.1) {\textbf{Degraded mel}};
\node[blk,text width=3.4cm] (bb) at (1.6,0) {\textbf{Shared MobileNetV3}\\backbone + signatures};
\node[los,text width=2.0cm] (lcls) at (5.4,1.35) {$\mathcal{L}_{\mathrm{cls}}$\\dual CE};
\node[los,text width=2.0cm] (lpc) at (5.4,0) {$\mathcal{L}_{\mathrm{PC}}$\\coherence};
\node[los,text width=2.0cm] (lft) at (5.4,-1.35) {$\mathcal{L}_{\mathrm{FTSC}}$\\signatures};
\draw[arr] (c)--(bb);
\draw[arr] (d)--(bb);
\draw[arr] (bb)--(lcls);
\draw[arr] (bb)--(lpc);
\draw[arr] (bb)--(lft);
\end{tikzpicture}}
\caption{FT-RCL paired clean/degraded-view training architecture.}
\label{fig:ft-rcl-paired}
\end{figure}


\subsection{Loss functions}
\label{sec:losses}

Let $z_c, z_d$ be the logits from the clean and degraded mel inputs, $q_c, q_d$ the corresponding fake-class probabilities, and $s_c, s_d$ the signatures.

\subsubsection{Dual classification}

With the class-weighted cross-entropy of Equation~\eqref{eq:ce},
\begin{equation}
\mathcal{L}_{\mathrm{cls}} = \mathcal{L}_{\mathrm{CE}}(z_c, y) + \beta\,\mathcal{L}_{\mathrm{CE}}(z_d, y), \qquad \beta = 0.5 .
\label{eq:cls}
\end{equation}

\subsubsection{Prediction coherence (PC)}

\begin{equation}
\mathcal{L}_{\mathrm{PC}} = \frac{1}{B}\sum_{i=1}^{B} \big(q_{c,i} - q_{d,i}\big)^2 .
\label{eq:pc}
\end{equation}
If, for one clip, $q_c=0.80$ and $q_d=0.55$, the clip adds $0.0625$ to the sum and pushes the degraded prediction toward the clean one. PC cannot tell right from wrong, since two views can agree on the wrong label. A squared difference of probabilities also has a weak gradient when one of the two outputs is saturated near 0 or 1, which is exactly the case of a confident clean prediction and a collapsed degraded one. A Kullback--Leibler or Jensen--Shannon term on the output distributions does not have this property~\cite{hendrycks2020augmix}; we did not compare them.

\subsubsection{Feature signature coherence (FTSC)}

For band $k$ and sample $i$, let $s_{c,i,k}, s_{d,i,k} \in \mathbb{R}^{C}$ be the L2-normalized band vectors and $\cos_{i,k} = s_{c,i,k}^\top s_{d,i,k}$. FTSC compares the clean and degraded signatures by the cosine similarity along the channel dimension for each band and averages over bands and batch:
\begin{equation}
\mathcal{L}_{\mathrm{FTSC}} = \frac{1}{B}\sum_{i=1}^{B} \frac{1}{K}\sum_{k=1}^{K} (1 - \cos_{i,k}).
\label{eq:ftsc}
\end{equation}
Figure~\ref{fig:pc-ftsc} illustrates PC and FTSC.

\figimage{pc-ftsc}{Prediction coherence and feature signature coherence.}{figures/plots/btp_report/pc_ftsc.png}

\subsubsection{Type, confidence, and severity robustness gate}

For sample $i$ with audio type $t_i$, the gate combines a type weight $\alpha_{t_i}$ (defaults: speech and singing $1.0$, music $0.3$, sound $0.5$), a confidence $\kappa_i = 2|q_{c,i} - 0.5|$ (treated as a constant, so no gradient flows through $\kappa$), and the severity factor $\eta(\rho_i)$ of Equation~\eqref{eq:eta}:
\begin{equation}
w_i = \alpha_{t_i}\,\kappa_i\,\eta(\rho_i), \qquad \bar{w} = \frac{1}{B}\sum_{i=1}^{B} w_i .
\label{eq:gate}
\end{equation}
The type weights were set by hand to reduce the constraint on the two types where the baseline is weakest; they were not estimated from data and no other values were tried. The confidence term switches the coherence loss off for uncertain clean predictions, which are the samples that probably need it most. Because the batch mean $\bar w$ multiplies the batch-level coherence loss, the gate scales the loss of the whole batch and does not weight individual samples. An ungated variant would show whether the gate matters; it was not run.

\subsection{Total loss and training modes}

The total loss of the full FT-RCL objective (mode M3) is
\begin{equation}
\mathcal{L}_{\mathrm{M3}} = \mathcal{L}_{\mathrm{cls}} + \bar{w}\,\Big(\gamma_{\mathrm{pc}}\mathcal{L}_{\mathrm{PC}} + \gamma_{\mathrm{ftsc}}\mathcal{L}_{\mathrm{FTSC}}\Big), \qquad \gamma_{\mathrm{pc}} = 1.0,\ \gamma_{\mathrm{ftsc}} = 0.5 .
\label{eq:m3}
\end{equation}
Table~\ref{tab:modes} summarizes how the four training modes combine the terms of Equations~\eqref{eq:cls}--\eqref{eq:ftsc} and the gate of Equation~\eqref{eq:gate}. Mode~M1 uses $\mathcal{L}_{\mathrm{CE}}(z_d,y)$ only; mode~M2 uses $\mathcal{L}_{\mathrm{cls}} + \bar{w}\,\gamma_{\mathrm{pc}}\mathcal{L}_{\mathrm{PC}}$ with $\eta \equiv 1$ and $\gamma_{\mathrm{ftsc}} = 0$.

\begin{table}[htbp]
\centering
\caption{Training modes and their loss components.}
\label{tab:modes}
\small
\renewcommand{\arraystretch}{1.2}
\begin{tabular}{@{}llll@{}}
\toprule
Mode & Network & Views used in the loss & Loss \\
\midrule
M0 & PANNs MobileNetV2 (E3) & Clean & $\mathcal{L}_{\mathrm{CE}}(z_c,y)$ \\
M1 & FT-RCL detector & Degraded & $\mathcal{L}_{\mathrm{CE}}(z_d,y)$ \\
M2 & FT-RCL detector & Clean and degraded & $\mathcal{L}_{\mathrm{cls}} + \bar{w}\,\gamma_{\mathrm{pc}}\mathcal{L}_{\mathrm{PC}}$ ($\eta\equiv 1$) \\
M3 & FT-RCL detector & Clean and degraded & $\mathcal{L}_{\mathrm{cls}} + \bar{w}\,(\gamma_{\mathrm{pc}}\mathcal{L}_{\mathrm{PC}} + \gamma_{\mathrm{ftsc}}\mathcal{L}_{\mathrm{FTSC}})$ \\
\bottomrule
\end{tabular}
\end{table}

\section{STAR-Edge extension}
\label{sec:star-edge-method}

STAR-Edge (Severity-Trajectory and Asymmetric Robustness for Edge) adds three terms to FT-RCL. It uses the same backbone, frontend, degradation catalogue and optimisation settings.

\paragraph{Asymmetric bona-fide robustness (ABR).} Genuine recordings suffer most when compression or noise creates artefacts that look like synthesis. For the real clips of a batch ($\mathcal{R}$), ABR penalises degraded views that look more fake than the clean view by more than a margin $m=0.05$:
\begin{equation}
\mathcal{L}_{\mathrm{ABR}} = \frac{1}{|\mathcal{R}|}\sum_{i\in\mathcal{R}} \max\big(0,\ q_{d,i}-q_{c,i}+m\big).
\label{eq:abr}
\end{equation}

\paragraph{Severity trajectory coherence (STC).} STC adds the signature and probability agreement between consecutive points of a trajectory $(s^{(0)},\ldots,s^{(T)})$:
\begin{equation}
\mathcal{L}_{\mathrm{STC}} = \sum_{t=0}^{T-1}\Big[\tfrac{1}{B}\sum_i\big(1-s^{(t)\top}_i s^{(t+1)}_i\big) + \tfrac{1}{B}\sum_i\big(q^{(t)}_i-q^{(t+1)}_i\big)^2\Big].
\label{eq:stc}
\end{equation}
In the implementation the trajectory has only two points, the clean and the degraded view, so STC reduces to a second copy of the pairwise terms of FT-RCL rather than a coherence across several severities. It is applied on a random half of the training steps.

\paragraph{Channel auxiliary detection (CAD).} Four waveform statistics (RMS, crest factor, spectral flatness, zero-crossing rate) are mapped by a small MLP to an eight-dimensional channel code $u$. The pooled embedding $z$ is projected orthogonally to $u$, $z'=z-(z^\top\hat u)\hat u$ with $\hat u=u/\lVert u\rVert$, and a linear head predicts the degradation family of the degraded clip from $u$ with a cross-entropy loss. The code is detached from the projection during the first three epochs.

With weights $\rho_{\mathrm{ABR}}=0.2$ and $\delta_{\mathrm{STC}}=0.3$, the objective is $\mathcal{L}_{\mathrm{M3}}+0.2\,\mathcal{L}_{\mathrm{ABR}}+\mathcal{L}_{\mathrm{CAD}}+0.3\,\mathcal{L}_{\mathrm{STC}}$. Like the other additions, these terms are only used in training. The signature problem described above applies here as well and is worse: with CAD the signature is computed from the projected pooled embedding, which has a single row, so only one of the eight bands was defined.

\section{Optimization and checkpoint selection}
\label{sec:optimization}

\subsection{Optimiser and schedule}

Training uses AdamW~\cite{loshchilov2019adamw} with learning rate $10^{-3}$ and weight decay $10^{-4}$, batch size 32, at most 20 epochs and early stopping with patience 5 on the selection score. The learning rate follows ReduceLROnPlateau on development Macro-F1 (factor 0.5, patience 2). The runs used a CPU, so no mixed precision was involved.

\subsection{Checkpoint selection}
\label{sec:selection}

The baseline keeps the epoch with the highest development Macro-F1. FT-RCL and STAR-Edge keep the epoch with the highest robustness-aware score
\begin{equation}
J = \mathrm{MacroF1}_{\mathrm{dev,clean}} - \lambda\,\widehat{\Delta \mathrm{F1}}_{\mathrm{dev}}, \qquad \lambda = 0.5 ,
\label{eq:jscore}
\end{equation}
where $\widehat{\Delta \mathrm{F1}}_{\mathrm{dev}}$ is the drop in Macro-F1 under two quick degradations. This makes the comparison unfair in a way that cannot be separated from the method: the models are chosen by different rules. The estimate itself was also unreliable. It used the first five batches of the development list, and the list is sorted by type and label, so these 160 clips were all fake music. Macro-F1 on a single class carries no information about robustness, and the logged $\widehat{\Delta \mathrm{F1}}$ values lie between $-0.03$ and $0.04$ without a trend. In addition the estimate can be negative, in which case $J$ rewards a model for scoring better on degraded than on clean audio. The corrected implementation uses a fixed development subset stratified by type and label, all degradation families and a drop clamped at zero, and the same selection rule is available for every model.

\section{Summary of implementation corrections}
\label{sec:audit}

Reviewing the code against the results exposed the problems below. The reported results of Chapter~\ref{ch:results} come from the implementation before the corrections, except where stated; the corrections are in the code base.

\begin{enumerate}
  \item The signature module produced empty bands, so the FTSC term was defined on three bands (one for STAR-Edge) and the logged loss was undefined.
  \item The severity curriculum did not reach the persistent data-loading workers.
  \item The estimate of the robustness drop used for checkpoint selection was computed on 160 fake music clips.
  \item Deterministic degradations used one noise seed for all clips, and a missing codec tool silently returned the clean waveform instead of failing.
  \item Checkpoints were loaded in a mode that ignores missing or unexpected weights.
  \item The robustness pilot used the first rows of a sorted list (Section~\ref{sec:robustness-pilot}).
\end{enumerate}

\section{Training versus inference}
\label{sec:train-infer}

\textbf{Training (M3)} requires paired degradations, two mel computations, dual logits and signatures, and the full loss $\mathcal{L}_{\mathrm{M3}}$. \textbf{Inference} applies $\phi$ once and performs a single forward pass; signatures and degradations are omitted. This preserves the latency expected of an edge-compatible deployment (Figure~\ref{fig:train-vs-infer}).

\begin{figure}[htbp]
\centering
\begin{tikzpicture}
\begin{scope}
\node[font=\scriptsize\bfseries\itshape,text=black!70,anchor=south west] at (-0.15,1.35) {Training (M3)};
\draw[dashed,rounded corners=4pt,draw=black!45] (-0.35,-0.85) rectangle (4.85,1.05);
\node[blk,text width=1.4cm] (tm) at (0,0.35) {\textbf{Mel}};
\node[blk,text width=1.4cm] (tc) at (2.2,0.35) {\textbf{CNN}};
\node[blk,text width=1.4cm] (ts) at (4.4,0.35) {\textbf{Loss}};
\draw[arr] (tm)--(tc);
\draw[arr] (tc)--(ts);
\node[font=\scriptsize\itshape,text=black!70,align=center,text width=4.6cm] at (2.2,-0.45) {Paired clean/degraded views\\dual CE + PC + FTSC + gate};
\end{scope}
\begin{scope}[xshift=6.5cm]
\node[font=\scriptsize\bfseries\itshape,text=black!70,anchor=south west] at (-0.15,1.35) {Inference (deploy)};
\draw[dashed,rounded corners=4pt,draw=deployedge] (-0.35,-0.85) rectangle (4.85,1.05);
\node[dep,text width=1.4cm] (im) at (0,0.35) {\textbf{Mel}};
\node[dep,text width=1.4cm] (ic) at (2.2,0.35) {\textbf{CNN}};
\node[dep,text width=1.4cm] (is) at (4.4,0.35) {\textbf{Score}};
\draw[arr] (im)--(ic);
\draw[arr] (ic)--(is);
\node[font=\scriptsize\itshape,text=black!70,align=center,text width=4.6cm] at (2.2,-0.45) {Single mel pass\\single forward pass};
\end{scope}
\end{tikzpicture}
\caption{Training versus inference workflow.}
\label{fig:train-vs-infer}
\end{figure}


\section{Evaluation metrics}
\label{sec:evaluation-metrics}

\subsection{Metrics}

For a clip with label $y \in \{0,1\}$, the model produces logits $z$ and probabilities $p = \mathrm{softmax}(z)$ (Section~\ref{sec:formulation}). The \textbf{detection score} is
\begin{equation}
q = p_1 = \frac{e^{z_1}}{e^{z_0} + e^{z_1}} \in [0,1],
\label{eq:q}
\end{equation}
the predicted probability that the clip is fake. For a decision threshold $\tau$, the predicted label is $\hat{y}_\tau = \mathbb{1}[q \ge \tau]$. Unless stated otherwise, $\tau = 0.5$. Let $\mathrm{TP}$, $\mathrm{FP}$, $\mathrm{TN}$, and $\mathrm{FN}$ denote the counts of true positives, false positives, true negatives, and false negatives, with the fake class ($y=1$) as the positive class. The metrics are defined as follows.

\paragraph{Accuracy:} the fraction of correct decisions at threshold $0.5$ on $q$,
\begin{equation}
\mathrm{Acc} = \frac{\mathrm{TP} + \mathrm{TN}}{\mathrm{TP} + \mathrm{TN} + \mathrm{FP} + \mathrm{FN}}.
\label{eq:acc}
\end{equation}

\paragraph{Macro-F1:} the unweighted mean of the F1 scores of the real and fake classes. For class $c \in \{0,1\}$, with precision $P_c$ and recall $R_c$,
\begin{equation}
\mathrm{F1}_c = \frac{2\,P_c R_c}{P_c + R_c}, \qquad
\mathrm{Macro\text{-}F1} = \tfrac{1}{2}\left(\mathrm{F1}_0 + \mathrm{F1}_1\right),
\label{eq:macrof1}
\end{equation}
where $P_1 = \mathrm{TP}/(\mathrm{TP}+\mathrm{FP})$, $R_1 = \mathrm{TP}/(\mathrm{TP}+\mathrm{FN})$, $P_0 = \mathrm{TN}/(\mathrm{TN}+\mathrm{FN})$, and $R_0 = \mathrm{TN}/(\mathrm{TN}+\mathrm{FP})$. Unless stated otherwise, the overall values are computed over all evaluation clips pooled together.

\paragraph{AUROC:} the area under the ROC curve traced by the score $q$. With the false-positive rate $\mathrm{FPR}(\tau) = \mathrm{FP}/(\mathrm{FP}+\mathrm{TN})$ and the true-positive rate $\mathrm{TPR}(\tau) = \mathrm{TP}/(\mathrm{TP}+\mathrm{FN})$ at threshold $\tau$,
\begin{equation}
\mathrm{AUROC} = \int_{0}^{1} \mathrm{TPR}\big(\mathrm{FPR}^{-1}(u)\big)\,du
= \Pr\big(q_{+} > q_{-}\big) + \tfrac{1}{2}\Pr\big(q_{+} = q_{-}\big),
\label{eq:auroc}
\end{equation}
where $q_{+}$ and $q_{-}$ are the scores of a randomly drawn fake and real clip, respectively.

\paragraph{EER:} the error rate at the operating point where the false-positive and false-negative rates are equal. With $\mathrm{FNR}(\tau) = \mathrm{FN}/(\mathrm{FN}+\mathrm{TP})$, it is computed at the ROC point where $\mathrm{FPR}$ and $\mathrm{FNR}$ are closest,
\begin{equation}
\tau^{\ast} = \arg\min_{\tau}\,\big|\mathrm{FPR}(\tau) - \mathrm{FNR}(\tau)\big|, \qquad
\mathrm{EER} = \tfrac{1}{2}\big(\mathrm{FPR}(\tau^{\ast}) + \mathrm{FNR}(\tau^{\ast})\big).
\label{eq:eer}
\end{equation}

\paragraph{AT-ADD Macro-F1:} for each audio type $t \in \mathcal{T} = \{\text{speech}, \text{singing}, \text{music}, \text{sound}\}$, the F1 scores of the real and fake classes are averaged, and the four type-wise scores are then averaged without weighting~\cite{xie2026atadd},
\begin{equation}
\mathrm{F1}_{\mathrm{AT\text{-}ADD}} = \frac{1}{|\mathcal{T}|}\sum_{t \in \mathcal{T}} \tfrac{1}{2}\left(\mathrm{F1}_{0}^{(t)} + \mathrm{F1}_{1}^{(t)}\right),
\label{eq:atadd}
\end{equation}
where $\mathrm{F1}_{c}^{(t)}$ is the F1 score of class $c$ computed on the clips of type $t$ only. Two further quantities are reported. Because a deployed detector has one threshold for all types, the per-type results are also given at a single global threshold, taken at the pooled EER point. Calibration is summarised by the expected calibration error over 15 equal-width bins~\cite{guo2017calibration}. Intervals are 95\% percentile intervals of a clip-level bootstrap with 1000 resamples~\cite{efron1993bootstrap}; differences between models use the same resamples, so they are paired. The robustness measure is the drop $\Delta M = M_{\mathrm{clean}}-M_{\mathrm{conditioned}}$; because Macro-F1 at a fixed threshold reacts strongly to a shift in the score distribution, $\Delta$AUROC is the more reliable of the three.

No settings were tuned on evaluation data.
